% Artigo: o t8Flow -- estrategias de um sistema de CFD sobre o t8code, e a
% verificacao da equacao de Poisson pelo metodo das solucoes manufaturadas.
%
% Compila com pdflatex + bibtex + pdflatex + pdflatex, ou latexmk -pdf main
% No Overleaf: enviar a pasta inteira; este e' o arquivo principal.
%
% MESMA CONVENCAO DOS OUTROS ARTIGOS DESTA ARVORE: os NUMEROS vivem em
% dados/*.dat, gerados por ./gera-dados.sh, e sao lidos por pgfplotstable --
% nunca digitados no texto.  Regenerar uma medida troca o .dat e o texto
% acompanha, sem a chance de a tabela e a prosa discordarem.
%
% O DESENHO DO ARTIGO e' uma matriz de ESQUEMA x MALHA, e a estrutura segue
% dela em vez de separar validacao de resultados:
%
%                       | malha uniforme | malha adaptada (quadtree 2:1)
%   fluxo 2a ordem      |   grau 2       |   grau 1
%   compacto 4a ordem   |   grau 3       |   recusa
%   MLS + Gauss         |   grau = ordem |   grau = ordem
%
% As nove celulas dessa matriz -- incluindo a recusa, que e' um resultado -- sao
% o que o artigo tem de trazer, e cada uma esta' medida.
%
% STATUS DE CADA SECAO -- o que tem respaldo medido e o que nao tem.  Cada secao
% abre com um bloco \status{} dizendo isso, e esses blocos somem do PDF quando
% \nostatus e' definido.  Existem para que ninguem -- nos inclusive -- leia um
% paragrafo de planejamento como resultado.
%
%   1 Introducao            PRONTA
%   2 Estrategias           PRONTA: todas as decisoes tomadas e implementadas
%   3 Aproximacoes          PRONTA: as cinco medidas, com os contrastes
%   4 Metodo de verificacao PRONTA
%   5 Resultados            PRONTA: malha uniforme e adaptada, 2D, serial
%   6 Verificacao da suite  PRONTA: 69 clausulas, 141 casos
%   7 Conclusoes            PRONTA
%
%   O QUE O ARTIGO NAO TRAZ, e diz onde: o CUSTO do paralelo, tres dimensoes, e
%   Robin dentro do operador.  Estao na secao 7.2 como trabalho declarado, nao
%   escondidos entre os resultados.  A CORRECAO do paralelo esta' na secao 8.

\documentclass[11pt,a4paper]{article}

\usepackage[T1]{fontenc}
\usepackage[utf8]{inputenc}
\usepackage[brazil]{babel}
\usepackage{lmodern}
\usepackage{microtype}
\usepackage{amsmath,amssymb}
\usepackage{booktabs}
\usepackage{multirow}
\usepackage{adjustbox}
\usepackage{siunitx}
\usepackage{graphicx}
\usepackage{pgfplots}
\usepackage{pgfplotstable}
\usepackage{tikz}
\usepackage{xcolor}
\usepackage[hidelinks]{hyperref}
\usepackage[margin=2.7cm]{geometry}

\pgfplotsset{compat=1.18}
\usetikzlibrary{positioning}

\newcommand{\dat}[1]{dados/#1}

% \status{...}: o bloco que diz qual respaldo a secao tem.  Definir \nostatus
% antes de \begin{document} para suprimi-los na versao de submissao.
\newcommand{\status}[1]{%
  \ifdefined\nostatus\else
    \begingroup
    \par\medskip\noindent
    \colorbox{yellow!18}{\begin{minipage}{0.97\linewidth}
      \small\textbf{Situação desta seção.} #1
    \end{minipage}}\par\medskip
    \endgroup
  \fi}

% \falta{...}: uma lacuna pontual dentro do texto, visivel de proposito.
\newcommand{\falta}[1]{\textcolor{red!70!black}{\textbf{[falta:} #1\textbf{]}}}
% \reserva: nao e' o mesmo que \falta.  \falta marca o que falta MEDIR;
% \reserva marca o que foi medido mas com uma ressalva que muda a leitura.
% Confundir os dois esconderia o segundo tipo entre pendencias, quando ele e'
% um resultado com reserva.
\newcommand{\reserva}[1]{\par\smallskip\noindent\textcolor{orange!60!black}{%
  \textbf{[reserva:} #1\textbf{]}}\par\smallskip}

% Tabela de convergencia lida de um .dat.  Onde o esquema e' EXATO o gerador
% escreve "---" na coluna de ordem, e a tabela o repassa: um esquema exato nao TEM ordem para medir, e
% imprimir um numero ali seria ruido dividido por ruido.
\newcommand{\tabconv}[2]{%
  \pgfplotstabletypeset[
    col sep=space,
    columns={nivel,celulas,h,einf,ordinf,eL2,ordL2},
    columns/nivel/.style={column name=nível, int detect},
    columns/celulas/.style={column name=células, int detect},
    columns/h/.style={column name=$h$, sci, sci zerofill, precision=3},
    columns/einf/.style={column name=$\|e\|_\infty$, sci, sci zerofill, precision=3},
    columns/ordinf/.style={column name=ordem, string type, column type=r},
    columns/eL2/.style={column name=$\|e\|_2$, sci, sci zerofill, precision=3},
    columns/ordL2/.style={column name=ordem, string type, column type=r},
    every head row/.style={before row=\toprule, after row=\midrule},
    every last row/.style={after row=\bottomrule},
  ]{\dat{#1}}%
  \par\smallskip\noindent\footnotesize #2\par}

\title{t8Flow: um sistema de escoamentos sobre florestas adaptativas\\
       \large Estratégias de projeto e verificação da equação de Poisson\\
       pelo método das soluções manufaturadas}
\author{\falta{autores}}
\date{12 de outubro de 2026}

\emergencystretch=3em
\begin{document}
\maketitle

\begin{abstract}
Este trabalho apresenta o \textbf{t8Flow}, um sistema de dinâmica de fluidos
computacional construído sobre a biblioteca de florestas adaptativas
\textsc{t8code}, e a verificação de seu operador de Poisson pelo método das
soluções manufaturadas. Duas decisões de projeto governam o resto: a
\emph{floresta é a estrutura} --- a malha não é produzida pelo t8code e copiada
para outra árvore, e nenhum processo detém a malha inteira --- e \emph{há um
domínio só}, de modo que nenhuma propriedade precise de tradução de índice entre
malhas paralelas. A segunda decisão vem de um defeito medido no sistema
antecessor, onde dois domínios coincidentes em malha uniforme tornavam um índice
errado acidentalmente certo, e a falha só aparecia com refino.

Reaproveita-se do antecessor~\cite{sousa2019} o que é independente de malha: o invólucro de
solvers lineares e o aproximador por mínimos quadrados móveis ponderados. Sobre
essa base verificam-se \textbf{cinco} aproximações para o laplaciano em volumes
finitos: balanço de fluxos de segunda ordem com fechamento de parede de três e de
dois pontos, o esquema compacto de quarta ordem em suas formas de Mehrstellen e
de derivada segunda implícita de Lele, e um esquema novo de reconstrução por
mínimos quadrados com quadratura de Gauss.

Os resultados estabelecem o grau de exatidão de cada esquema em malha uniforme e
em malha \emph{quadtree} com transições 2:1, e três achados estruturam a
discussão. Primeiro, a exatidão do grau~2 do esquema de fluxo é propriedade da
malha uniforme e se perde na adaptada, onde ele é exato apenas para campo linear.
Segundo, o esquema compacto atinge quarta ordem mas é exato só até o grau~3, e o
limite é o fechamento de parede, não o interior. Terceiro, a aproximação por
mínimos quadrados com quadratura de Gauss é exata até o grau da ordem do ajuste
\emph{tanto em malha uniforme como em adaptada}, e o fator que a destravava não
era o fluxo na face: era a quadratura do termo de fonte no volume, cuja regra do
ponto médio mantinha o esquema inteiro em segunda ordem.

A execução paralela é tratada em seção própria, e por um motivo metodológico:
tudo nela é invisível em execução serial, e três defeitos reais do sistema viviam
exatamente nessa invisibilidade --- entre eles um que fazia um esquema
\emph{exato} para campo linear errar \num{1.70e-1} com dois processos, onde com
um errava \num{3.97e-13}.

Toda afirmação de exatidão vem acompanhada de um contraste --- a mesma medida
onde ela deve falhar --- e a suíte de testes nasceu junto com o código, com
69 cláusulas de contrato lidas pelo próprio arreio de execução, que exige
um caso passando para cada uma.
\end{abstract}

\tableofcontents

\section{Introdução}
\label{sec:introducao}

\status{Pronta. O contexto e as decisões aqui estão todos implementados e
cobertos por cláusulas do contrato de testes; as referências ao sistema
antecessor vêm do código dele, não de memória.}

Sistemas de dinâmica de fluidos computacional com refino adaptativo de malha
enfrentam, cedo ou tarde, a mesma pergunta de projeto: \emph{quem é dono da
estrutura da malha}. A resposta determina o que escala, o que pode ser refinado,
e onde um erro de índice se esconde.

O sistema antecessor deste trabalho --- o \textsc{HiGFlow}, cujo método de
diferenças finitas com interpolação sem malha em árvores não graduadas está em
Sousa \emph{et al.}~\cite{sousa2019} --- resolveu essa pergunta de um modo que
funcionou por anos e esgotou-se: a árvore hierárquica própria é dona da
estrutura, e uma biblioteca externa de florestas
adaptativas, a \textsc{t8code}~\cite{t8code}, entra apenas como \emph{produtora}
de malha. O levantamento que precedeu este trabalho mediu o limite dessa escolha
no próprio código: a fonte que produz malha pelo t8code declara, em letras, que
só sabe tratar malha uniforme, um nível e um \emph{patch} por bloco, e
\emph{recusa} refino em vez de aproximar.

Essa é uma recusa honesta, e é também o fim da linha. Um sistema que queira as
propriedades de escalonamento do t8code precisa do inverso: a floresta sendo a
estrutura, com o resto lendo dela. Esse é o t8Flow.

\subsection{O que este artigo traz}

Duas coisas, e elas se sustentam mutuamente.

A primeira é o conjunto de \textbf{estratégias de projeto} (seção
\ref{sec:estrategias}): quais decisões foram tomadas, quais alternativas foram
medidas antes de descartar, e --- no caso mais importante --- qual defeito
concreto do antecessor motivou qual decisão. Decisões de arquitetura justificadas
por gosto envelhecem mal; as daqui estão justificadas por medida ou por defeito
observado.

A segunda é a \textbf{verificação do operador de Poisson} pelo método das
soluções manufaturadas (seções \ref{sec:metodo} e \ref{sec:resultados}). A
equação de Poisson é o lugar certo para começar: ela é o passo de projeção de
qualquer método de passo fracionado para escoamento incompressível, é o passo
mais caro, e é onde uma malha adaptativa cobra o preço da transição entre níveis.
Um operador de Poisson que não seja exato para os polinômios que promete não
produzirá ordem em escoamento algum.

\subsection{Cinco aproximações, não uma}

O ponto central do artigo é que \textbf{cinco} aproximações distintas do
laplaciano foram implementadas e medidas no \emph{mesmo} arreio --- mesma malha,
mesmo solver linear, mesmas condições de contorno, mesmo catálogo de soluções
manufaturadas --- de modo que a diferença medida seja atribuível ao esquema e não
ao arranjo:

\begin{enumerate}
\item balanço de fluxos de segunda ordem com fechamento de parede de
      \textbf{três pontos};
\item o mesmo com fechamento de \textbf{dois pontos}, que serve de contraste;
\item o compacto de quarta ordem na forma \textbf{Mehrstellen} (nove pontos, com
      correção $h^2\nabla^2 f/12$ no termo de fonte);
\item o compacto de quarta ordem na forma de \textbf{derivada segunda implícita}
      de Lele~\cite{lele1992};
\item reconstrução por \textbf{mínimos quadrados móveis com quadratura de
      Gauss}, apresentada aqui.
\end{enumerate}

O quinto é novo neste sistema e existe por uma razão que os resultados tornam
precisa: é o único dos cinco que passa de segunda ordem \emph{em malha
adaptada}. Os dois compactos não servem em malha adaptada por construção --- o
estêncil de nove pontos pressupõe espaçamento uniforme --- e o de fluxo perde,
na transição 2:1, a exatidão de grau~2 que tinha na malha uniforme.

\subsection{A disciplina de oráculo}

Uma observação metodológica, porque ela governa a forma de toda a seção de
resultados. Neste projeto, \emph{nenhuma afirmação de exatidão aparece sozinha}.
Cada uma vem com um contraste: a mesma medida feita onde ela \textbf{deve}
falhar. Um teste que afirma ``o erro é de máquina'' sem exibir a condição em que
o erro deixa de ser de máquina não distingue um esquema correto de um oráculo
cego --- e a seção \ref{sec:verificacao} relata quatro defeitos reais deste
trabalho que \emph{apenas} o contraste encontrou, incluindo um que produzia erro
de \num{3.3e-1} em um nível de refino cercado por \num{1e-10} nos vizinhos.

\section{Estratégias de projeto}
\label{sec:estrategias}

\status{Pronta. Todas as decisões descritas estão implementadas; as que foram
tomadas após medida trazem a medida. A seção \ref{sec:est-escala} traz números de
escalabilidade em duas dimensões e a seção \ref{sec:par-escala} os de três,
todos medidos nesta máquina.}

\subsection{A floresta é a estrutura}

A decisão da qual as outras dependem. No t8Flow a malha \textbf{é} a floresta do
t8code: a topologia, a numeração, os vizinhos e a partição vêm dela, e nenhuma
estrutura paralela as duplica. Isso tem três consequências imediatas:

\begin{itemize}
\item \textbf{nenhum processo detém a malha inteira.} A curva de preenchimento
      de espaço (ordem de Morton) particiona as folhas, e cada processo conhece
      as suas mais uma camada de fantasmas de face;
\item \textbf{o refino é da floresta.} Adaptar é chamar o adaptador do t8code e
      \emph{reconfirmar} o balanceamento 2:1; não há sincronização entre duas
      estruturas que possam divergir;
\item \textbf{a partição é do t8code.} Não há cálculo próprio de balanceamento de
      carga. Essa é uma renúncia deliberada, e a seção \ref{sec:est-escala} a
      quantifica.
\end{itemize}

\subsection{Multiblocos é multiárvore}

O t8Flow é multiblocos, como o antecessor, e essa foi a peça que encaixou sem
atrito: a malha grossa (\emph{cmesh}) do t8code é nativamente multiárvore, e
\textbf{um bloco é uma árvore}. Blocos vizinhos são unidos pelas faces
coincidentes; uma sobreposição parcial de faces é recusada com exceção, porque
unir faces que não coincidem exatamente produziria uma topologia plausível e
errada.

\subsection{Um domínio só}
\label{sec:est-dominio}

Esta decisão vem de um defeito concreto do antecessor, e por isso é relatada com
ele.

No antecessor, o domínio da fração volumétrica do método VOF é, por construção,
uma malha \emph{uniforme no nível mais fino}, separada da árvore adaptativa que
carrega a quantidade de movimento. Sete funções de advecção tridimensional
percorriam a árvore adaptativa e usavam o índice obtido ali \emph{diretamente}
como índice de uma propriedade que vive no outro domínio. As contagens medidas
nos campos salvos:

\begin{center}
\begin{tabular}{lr}
\toprule
propriedade & células \\
\midrule
fração volumétrica (domínio uniforme) & \num{375000} \\
velocidade (árvore adaptativa)        & $\approx$ \num{81000} \\
\bottomrule
\end{tabular}
\end{center}

\textbf{Em malha uniforme os dois domínios coincidem célula a célula, e o índice
errado é acidentalmente o certo.} Foi por isso que todo o conjunto de testes
tridimensionais passou por anos. Com um nível de refino, oitenta mil iterações
escreviam em um vetor de trezentos e setenta e cinco mil posições sem relação
geométrica com a célula de origem --- e a fração parou de advectar, em silêncio,
enquanto a quantidade de movimento continuava correta.

A decisão para o sistema novo é, portanto, \textbf{um domínio só}: toda
propriedade vive na mesma floresta. Dois domínios paralelos exigiriam tradução de
índice explícita e verificada, e herdar o arranjo sem a tradução herdaria a
classe inteira de defeito. Paga-se barato no início e caro depois.

\subsection{Herança do que é independente de malha}

Duas peças do antecessor~\cite{sousa2019} foram reaproveitadas sem modificação, depois de
verificado que não dependem da estrutura de malha:

\begin{description}
\item[o invólucro de solvers lineares,] com implementações para PETSc,
      \emph{hypre}, SOR, um solver centralizado e um que escreve o sistema em
      disco. A interface recebe linhas, um lado direito, e devolve a solução;
      nada nela conhece malha;
\item[o aproximador por mínimos quadrados móveis ponderados] (seção
      \ref{sec:aprox-mls}) --- a interpolação sem malha que é a contribuição
      central de~\cite{sousa2019} e o que permite àquele método operar em
      árvores \emph{não graduadas} ---, que recebe uma nuvem de pontos e um
      ponto-alvo e devolve pesos. Também não conhece malha.
\end{description}

Uma verificação metodológica vale registro, porque a primeira tentativa dela foi
\emph{inválida}: para decidir se um arquivo é independente de malha, tentou-se
compilá-lo com um diretório de cabeçalhos-sombra vazios. O teste não provou nada,
porque a diretiva \verb|#include "x.h"| procura primeiro o diretório do próprio
arquivo-fonte --- de modo que a sombra nunca se aplicava, e um arquivo que
\emph{depende} de malha foi declarado independente. O conserto foi copiar cada
fonte para fora da árvore antes de compilar.

\subsection{Duas e três dimensões em bases separadas}

O t8code trata malhas bidimensionais e tridimensionais; o t8Flow mantém
\textbf{bases de código separadas} para cada uma, começando pela
bidimensional. A alternativa --- uma base só, com a dimensão como parâmetro de
compilação --- economiza linhas e cobra em outro lugar: cada laço sobre direções,
cada contagem de vizinhos de face e cada fórmula de quadratura ganha um caso que
não é exercitado no regime em que se está trabalhando. Com bases separadas, o
código bidimensional é lido e testado como código bidimensional.

\subsection{Armazenamento: centros e facetas}

As propriedades escalares e tensoriais vivem nos \textbf{centros} das células; as
vetoriais, nas \textbf{facetas} --- o arranjo deslocado (MAC), que é o que torna
o acoplamento pressão--velocidade estável sem estabilização adicional.

A numeração de facetas exige uma decisão de propriedade, e ela é: a faceta
pertence à célula para a qual é a face \emph{``menos''} em cada direção, e
exatamente a uma. As faces ``mais'' que caem no contorno do domínio recebem
faceta própria, porque não há célula à direita que as possua. Em uma malha
uniforme de $N\times N$ isso dá $N(N+1)$ facetas por direção, e essa contagem é
afirmada pela suíte --- é barata e pega erro de propriedade antes que ele apareça
como divergência que não fecha.

\subsection{Estênceis em lote, e a GPU quando houver}
\label{sec:est-gpu}

O estêncil de ordem arbitrária é pedido ponto a ponto ou \textbf{em lote}: $N$
pontos de consulta, $N$ estênceis, em uma chamada. A forma de pedir foi definida
antes do \emph{back-end}, porque é ela que decide se a execução em acelerador é
possível.

\begin{description}
\item[seleção do suporte] fica sempre na CPU --- é navegação de malha, busca por
      vizinhança, crescimento por anéis;
\item[a resolução dos sistemas normais] é o que vai para o acelerador, com
      implementação em \textbf{OpenCL} (escolhido por portabilidade: funciona com
      placas AMD via \emph{rusticl}, Intel via NEO, e NVIDIA);
\item[o \emph{back-end} de CPU existe sempre,] e a escolha entre os dois é
      \textbf{sinalizador de compilação}, não detecção silenciosa em tempo de
      execução.
\end{description}

Precisão simples basta para o estêncil, e isso foi uma decisão consciente: os
pesos entram depois em um sistema linear resolvido em precisão dupla, e o erro
dos pesos é de truncamento do ajuste, não de arredondamento.

O oráculo desse caminho é \emph{ponto a ponto contra a CPU}, e ele encontrou um
defeito que a reprodução polinomial não encontrava: o núcleo em OpenCL usava
mínimos quadrados \textbf{não ponderados}, enquanto o núcleo herdado usa o peso
móvel $1/(d^2+\varepsilon)$. Ambos reproduzem polinômios do grau da ordem
exatamente --- e por isso o teste de reprodução passava nos dois ---, mas com
\emph{pesos diferentes}: diferença medida de \num{0.311} em um peso individual.
Ponderado corretamente, a diferença caiu para \num{8.8e-11}.

Dois percalços de medição no mesmo caminho, registrados porque são o tipo de erro
que faz um relatório afirmar o contrário do que mediu. Primeiro, recompilar com o
sinalizador trocado \emph{não} reconstruía, porque o sinalizador não estava nas
dependências do \emph{make}; o conserto foi gravar a configuração em um arquivo
conferido na leitura do \emph{makefile}. Segundo, a diferença GPU--CPU medida deu
\textbf{exatamente zero} --- porque com a GPU compilada ela é o \emph{back-end}
padrão, e a ``referência de CPU'' estava sendo calculada na GPU. O caso agora
força a CPU explicitamente e \emph{rejeita} o zero exato como suspeito.

\subsection{A escalabilidade, medida}
\label{sec:est-escala}

Duas medidas, nesta máquina, em duas dimensões.

A primeira é o \textbf{ganho de divisão}: \num{9377} células em quatro processos
renderam apenas \num{1.6}$\times$. A leitura não é que não escala --- é que
\num{2300} células por processo é pouco trabalho para o custo de comunicação. O
regime em que a divisão vence é o de mais células por processo, e ele é limitado
por \emph{tamanho}, não por dimensão.

O regime tridimensional, onde o antecessor tem \num{375000} células --- quarenta
vezes mais trabalho por processo ---, passou a ser medido e está na seção
\ref{sec:par-escala}: \num{4.92}$\times$ em oito processos no esquema de fluxo,
com \num{262144} células e erro idêntico ao serial.

\reserva{essa medida foi feita em \textbf{oito núcleos físicos}, e é o teto
desta máquina. O escalonamento além disso --- que é onde o regime do antecessor
de fato vive --- continua sendo expectativa, não resultado.}

A segunda é o \textbf{desbalanceamento por refino local}: refinar aleatoriamente
desbalanceia a partição em \num{2.8}$\times$, e reparticionar corrige isso a um
custo de 3 a 6\% do tempo de passo.

\subsection{A suíte nasce com o código}

A última estratégia é de processo, e foi pedida explicitamente: a suíte de testes
nasce junto com o código, e não depois dele. A forma que isso toma aqui é um
\textbf{contrato} --- um cabeçalho com cláusulas numeradas, cada uma declarando
uma propriedade que o sistema afirma e nomeando o caso que a verifica. O arreio
de execução \emph{lê as cláusulas do contrato} e exige um caso passando para cada
uma: cláusula sem teste sai como falha, não como omissão. A seção
\ref{sec:verificacao} descreve o estado e o que ele já pegou.

\section{As aproximações}
\label{sec:aproximacoes}

\status{Pronta. As cinco aproximações estão implementadas e medidas no mesmo
arreio; os coeficientes abaixo são os do código, não os da literatura reescritos.
A equivalência entre as duas formas do compacto foi verificada, não suposta
(seção \ref{sec:aprox-equiv}).}

Esta seção descreve o aproximador por mínimos quadrados móveis, que é a base
comum, e em seguida as cinco aproximações do laplaciano que foram medidas.

\subsection{Mínimos quadrados móveis ponderados}
\label{sec:aprox-mls}

O aproximador é o da família de superfícies por mínimos quadrados
móveis~\cite{lancaster1981}, na forma em que o sistema antecessor o emprega para
interpolar em árvores não graduadas~\cite{sousa2019}.

Dado um ponto-alvo $\mathbf{x}$ e uma nuvem de suporte $\{\mathbf{x}_k\}_{k=1}^n$
com valores $u_k$, o ajuste procura o polinômio $p$ de grau $m$ que minimiza

\begin{equation}
  \sum_{k=1}^{n} D_k^2 \bigl[p(\mathbf{x}_k) - u_k\bigr]^2,
  \qquad
  D_k^2 = \frac{1}{\|\mathbf{x}_k - \mathbf{x}\|^2 + \varepsilon},
  \label{eq:mls-funcional}
\end{equation}

e devolve, não o polinômio, mas o vetor de \textbf{pesos} $w_k$ tal que o
funcional desejado de $p$ em $\mathbf{x}$ seja $\sum_k w_k u_k$. O peso móvel de
\eqref{eq:mls-funcional} --- cada ponto pesado pelo inverso do quadrado da
distância --- é o que faz o ajuste ser \emph{móvel}: o polinômio que sai depende
do ponto onde se pergunta, e não é um ajuste global.

\paragraph{A base é relativa e adimensionalizada.} Os monômios são escritos na
variável deslocada e escalada
\begin{equation}
  \mathbf{s}_k = \frac{\mathbf{x}_k - \mathbf{x}}{L},
  \qquad
  L = \max_k \|\mathbf{x}_k - \mathbf{x}\|,
  \label{eq:mls-escala}
\end{equation}
e em duas dimensões a base tem $3$, $6$, $10$ e $15$ monômios para as ordens
$1$ a $4$. O deslocamento é o que torna o funcional de avaliação trivial: no
alvo, $\mathbf{s}=\mathbf{0}$, e portanto o valor é o coeficiente do monômio
constante, $\partial/\partial x$ é o coeficiente de $s_x$, e assim por diante.

A \emph{escala} de \eqref{eq:mls-escala} é menos óbvia e foi necessária. Sem ela,
a matriz normal da ordem~4 tem número de condição da ordem de $h^{-8}$; com
$h=\num{3e-2}$ isso passa de $10^{12}$, e o ajuste degrada \textbf{em silêncio}
--- devolvendo pesos, não erro. Com a escala, o condicionamento não depende de
$h$, e a cláusula D2 do contrato afirma isso pelo caminho em que falharia: a
mesma reprodução de grau~4 em um suporte mil vezes menor ($h=\num{1e-3}$).

O peso $D_k^2$, note-se, \textbf{permanece na escala física}. Ele é o peso do
núcleo herdado, e alterá-lo mudaria o ajuste em vez de apenas recondicioná-lo ---
o que tornaria as duas implementações incomparáveis. Essa comparabilidade é o que
permite a cláusula N7: os pesos de valor da implementação própria concordam com
os do núcleo herdado, que resolve o mesmo problema por decomposição QR, com
diferença máxima de \num{1.249e-16}. O caso \emph{rejeita} a concordância
exatamente nula como suspeita.

\paragraph{Resolução e os três funcionais.} Com $b_i$ os monômios, a matriz
normal e o lado direito são
\begin{equation}
  M_{ij} = \sum_k D_k^2\, b_i(\mathbf{s}_k)\, b_j(\mathbf{s}_k),
  \qquad
  M\mathbf{y} = \mathbf{e}_q,
  \qquad
  w_k = D_k^2 \sum_i y_i\, b_i(\mathbf{s}_k),
  \label{eq:mls-normal}
\end{equation}
onde $q$ seleciona o funcional: $q=0$ dá o \textbf{valor}, $q=1$ e $q=2$ dão
$\partial/\partial x$ e $\partial/\partial y$. Para os funcionais de derivada os
pesos são divididos por $L$, desfazendo a escala da variável.

Duas conferências baratas e independentes do campo acompanham cada ajuste: os
pesos de \emph{valor} somam $1$ (um ajuste que reproduz a constante não pode
fazer outra coisa) e os de \emph{derivada} somam $0$. São conferências que pegam
erro de escala e de sinal que a reprodução de um campo específico pode absorver.

\paragraph{O rebaixamento de ordem é registrado, nunca silencioso.} Quando o
suporte não sustenta a ordem pedida, o ajuste cai para a ordem que ele sustenta,
e um contador registra. A decisão é tomada pela razão entre o menor e o maior
pivô da eliminação --- não é o número de condição, mas é barata, monótona nele, e
a faixa é inequívoca: um suporte saudável dá razão entre $10^{-2}$ e $10^{-5}$, e
os ajustes que produziam o defeito descrito na seção \ref{sec:verificacao} ficavam
abaixo de $10^{-13}$. O limiar está na faixa vazia entre os dois regimes.

\paragraph{O estêncil: suporte por anéis e fechamento de contorno.} O suporte
cresce por \textbf{anéis} de vizinhos de face a partir da célula que contém o
ponto, e a parada só é avaliada com o anel \emph{fechado por inteiro} --- parar no
meio de um anel faria o suporte, e portanto os pesos, depender da ordem de
visita. Células de contorno entram pela relação afim descrita na seção
\ref{sec:aprox-bc}, de modo que o valor em cada ponto de suporte é
$a\,u[\mathrm{id}] + b$: o termo $a$ vai para a matriz e o termo $b$ para o lado
direito. Um único laço trata os três tipos de condição, e reusá-lo entre o
funcional de valor e o de derivada não é economia --- é o que garante que o
fechamento entre nos dois do mesmo modo.

\subsection{Balanço de fluxos de segunda ordem}
\label{sec:aprox-fluxo}

A aproximação de referência, em volumes finitos:
\begin{equation}
  \int_{\Omega_c} \nabla^2 u \,\mathrm{d}A
  \;=\; \oint_{\partial\Omega_c} \frac{\partial u}{\partial n} \,\mathrm{d}s
  \;\approx\; \sum_{f} \frac{|f|}{d_f}\,\bigl(u_{\mathrm{viz}} - u_c\bigr),
  \label{eq:fv}
\end{equation}
com $|f|$ a medida da faceta e $d_f$ a distância entre os centros. Na transição
2:1 os três fatores mudam juntos, e a implementação os escreve de forma única em
termos de $\varphi = 2^{\,n_c - n_{\mathrm{viz}}}$:
\begin{equation}
  h_{\mathrm{viz}} = \varphi\,h_c,
  \qquad
  |f| = \ell\,\min(1,\varphi),
  \qquad
  d_f = \tfrac{1}{2}\bigl(h_c + h_{\mathrm{viz}}\bigr),
  \label{eq:fv-transicao}
\end{equation}
mais uma \textbf{correção transversal}, porque na transição o segmento que une os
dois centros não é normal à faceta. A correção usa o gradiente transversal tomado
na célula \emph{grossa}, e seu sinal depende apenas do deslocamento transversal
entre os centros --- não de qual dos dois lados é o grosso. Fazê-lo depender disso
\textbf{dobrou} o erro na medição (\num{7.8e-2} contra \num{3.9e-2}), e esse foi o
sintoma que identificou a fórmula errada.

\subsection{O fechamento de parede}
\label{sec:aprox-bc}

\paragraph{A relação afim.} Toda condição de contorno entra pela mesma forma
$\alpha u + \beta\,\partial u/\partial n = g$ na parede, convertida em
$u_{\text{fantasma}} = A\,u_c + B$ com
\begin{equation}
  A = \frac{\beta - \alpha h/2}{\alpha h/2 + \beta},
  \qquad
  B = \frac{g\,h}{\alpha h/2 + \beta}.
  \label{eq:bc-afim}
\end{equation}
Os dois casos limites saem de \eqref{eq:bc-afim} sem caso especial: Dirichlet
($\beta=0$) dá $A=-1$, $B=2g$ --- reflexão em torno do valor na parede; Neumann
($\alpha=0$) dá $A=1$, $B=hg$ --- extrapolação pela derivada. Robin é o caso
geral. A escolha de uma relação afim única, em vez de três ramos, é o que permite
que o mesmo laço de montagem sirva aos três.

\paragraph{Três pontos, e o contraste de dois.} No esquema de fluxo, a derivada
normal na parede é construída com a célula de contorno, a vizinha e a seguinte,
a distâncias $s_1=h/2$ e $s_2 = s_1 + (h+h_{\mathrm{viz}})/2$ da parede:
\begin{equation}
  \frac{\partial u}{\partial n}\bigg|_{\text{parede}}
  \approx
  \underbrace{\frac{-s_1-s_2}{s_1 s_2}}_{c_g} u_g
  + \underbrace{\frac{s_2}{s_1(s_2-s_1)}}_{c_1} u_1
  + \underbrace{\frac{-s_1}{s_2(s_2-s_1)}}_{c_3} u_3 .
  \label{eq:fech3}
\end{equation}
O fechamento de \emph{dois} pontos permanece disponível e serve de contraste
controlado: com ele, o grau~2 deixa de ser exato e passa a convergir em segunda
ordem --- verde que nunca foi vermelho não prova nada.

Um erro de medida vale registro, porque foi instrutivo: em uma versão, $s_2$ foi
medido a partir do \emph{centro} da célula em vez da parede, de modo que na malha
uniforme ele valia $h$ em vez de $3h/2$. A regressão de grau~2 acusou
instantaneamente --- \num{1.97e-1} contra \num{3e-13}.

\subsection{Compacto de quarta ordem: Mehrstellen}
\label{sec:aprox-mehrstellen}

Na malha uniforme, o estêncil compacto de nove pontos:
\begin{equation}
  \frac{1}{6h^2}
  \begin{bmatrix}
    1 & 4 & 1\\
    4 & -20 & 4\\
    1 & 4 & 1
  \end{bmatrix} u
  \;=\;
  f + \frac{h^2}{12}\nabla^2 f .
  \label{eq:mehrstellen}
\end{equation}
A correção $h^2\nabla^2 f/12$ no lado direito é o que compra a quarta ordem a
partir de um estêncil que, sem ela, é de segunda.

O fechamento de parede aqui é um \textbf{fantasma de quatro pontos}, obtido da
interpolação que reproduz o valor prescrito na parede:
\begin{equation}
  u_{\text{fantasma}} = \tfrac{16}{5}g - 3u_0 + u_1 - \tfrac{1}{5}u_2 ,
  \label{eq:fantasma4}
\end{equation}
aplicado recursivamente, com os cantos tratados por composição. E é esse
fechamento, não o interior, que limita a \textbf{exatidão} a grau~3 --- o que a
seção \ref{sec:resultados} mede.

Este esquema \textbf{recusa} o que não sabe tratar: malha com mais de um bloco, e
malha adaptada. O estêncil de nove pontos pressupõe espaçamento uniforme, e
aplicá-lo em uma transição 2:1 daria um número plausível e errado.

\subsection{Compacto de quarta ordem: derivada segunda implícita}
\label{sec:aprox-lele}

A segunda forma do compacto aproxima a derivada segunda \emph{implicitamente},
na família de Lele~\cite{lele1992}:
\begin{equation}
  \tfrac{1}{10}u''_{i-1} + u''_i + \tfrac{1}{10}u''_{i+1}
  \;=\;
  \frac{6}{5}\,\frac{u_{i-1}-2u_i+u_{i+1}}{h^2}.
  \label{eq:lele}
\end{equation}
Em duas dimensões o operador aplicado é o produto tensorial $A_x A_y$ dos
operadores implícitos de cada direção, e o lado direito é $A_x A_y f$ dividido
pelo fator $(6/5)^2 = \num{1.44}$.

\subsection{As duas formas do compacto são o mesmo esquema}
\label{sec:aprox-equiv}

Isso foi \textbf{verificado, não suposto}: as duas formas produzem resultados
idênticos para soluções polinomiais de grau até~4 e diferem em \num{1.8}\% no
caso trigonométrico, o que é compatível com a mesma ordem de truncamento e
constantes distintas. Afirmar a equivalência sem medi-la seria afirmar o que a
álgebra sugere mas a implementação pode contradizer --- e o fator \num{1.44} é
exatamente o lugar onde uma implementação contradiz a álgebra em silêncio.

\subsection{Reconstrução por mínimos quadrados com quadratura de Gauss}
\label{sec:aprox-mlsgauss}

A quinta aproximação é a que fecha o vão entre as outras quatro: as duas formas
do compacto não servem em malha adaptada, e o de fluxo perde nela a exatidão de
grau~2. O fluxo pela faceta é
$\int_f \partial u/\partial n \,\mathrm{d}s$, e o esquema leva \textbf{os dois
fatores} à ordem alta, porque levantar um só não passa de segunda ordem:

\begin{center}
\begin{tabular}{lll}
\toprule
fator & fluxo de segunda ordem & reconstrução MLS \\
\midrule
derivada no ponto & diferença de dois centros & ajuste MLS de ordem 3 ou 4 \\
quadratura da faceta & ponto médio (exata p/ linear) & Gauss de 2 pontos (exata p/ cúbico) \\
quadratura do volume & $f(c)\,A$ (exata p/ linear) & Gauss $2\times2$ (exata p/ bicúbico) \\
\bottomrule
\end{tabular}
\end{center}

A terceira linha é o achado central deste trabalho, e a seção
\ref{sec:res-fonte} a mede isoladamente.

\paragraph{A quadratura da faceta.} Com $\mathbf{c}_f$ o centro da sub-faceta e
$\ell$ a sua medida, os dois nós de Gauss e os pesos normalizados para média são
\begin{equation}
  \mathbf{x}_q = \mathbf{c}_f \pm \frac{1}{\sqrt{3}}\,\frac{\ell}{2}\,\hat{\mathbf{t}},
  \qquad w_q = \tfrac{1}{2},
  \qquad
  \text{fluxo} = \ell \sum_{q=1}^{2} w_q \left.\frac{\partial u}{\partial n}\right|_{\mathbf{x}_q} .
  \label{eq:gauss-faceta}
\end{equation}
Na transição 2:1 a faceta da célula grossa é coberta por \textbf{duas}
sub-facetas, cada uma do tamanho da face da vizinha fina, e o fluxo total é a
soma. Integrar a face inteira de uma vez daria o fluxo médio sobre ela, que não é
a soma dos dois fluxos finos quando a derivada varia ao longo dela.

\paragraph{A quadratura do volume.} A fonte é integrada por Gauss $2\times2$,
\begin{equation}
  \int_{\Omega_c} f \,\mathrm{d}A \approx
  A \sum_{q_x=1}^{2}\sum_{q_y=1}^{2} w_{q_x} w_{q_y}\, f(\mathbf{x}_{q_x q_y}),
  \label{eq:gauss-volume}
\end{equation}
exata para $f$ bicúbico, o que cobre $f=\nabla^2 u$ de qualquer $u$ de grau até~5.

\paragraph{Conservação, por construção.} O suporte do estêncil de derivada
depende apenas do \textbf{ponto} de quadratura, nunca da célula que pergunta.
Logo as duas células que compartilham uma faceta obtêm o mesmo fluxo, sem que
ninguém combine nada --- inclusive na transição 2:1, onde a grossa parte a faceta
e a fina não parte a sua. Isso importa porque um esquema \textbf{não conservativo
ainda converge}: nenhum dos testes de ordem de convergência pegaria um erro de
indexação na sub-faceta. A cláusula D7 afirma a propriedade pelo caminho em que
ela falharia, com o oráculo refazendo a geometria das sub-facetas por conta
própria e comparando os estênceis peso a peso.

\section{Método de verificação}
\label{sec:metodo}

\status{Pronta. O catálogo de soluções, as normas e o oráculo de resíduo do
operador estão todos implementados e são o que produziu os números da seção
\ref{sec:resultados}.}

\subsection{O problema e as soluções manufaturadas}

Resolve-se, no quadrado unitário $\Omega=[0,1]^2$,
\begin{equation}
  \nabla^2 u = f \quad\text{em } \Omega,
  \qquad
  \alpha u + \beta\,\frac{\partial u}{\partial n} = g \quad\text{em } \partial\Omega .
  \label{eq:poisson}
\end{equation}

O método das soluções manufaturadas inverte a direção usual: escolhe-se $u$,
\emph{deriva-se} $f=\nabla^2 u$ analiticamente, impõem-se as condições de
contorno a partir do próprio $u$, e compara-se a solução discreta com $u$. A
vantagem decisiva é que a solução exata é conhecida em todo ponto, e portanto o
erro é mensurável sem refino de referência.

O catálogo tem seis soluções, com todas as derivadas escritas a mão --- nunca por
diferenciação numérica, que introduziria o erro que se quer medir:

\begin{center}
\begin{tabular}{lll}
\toprule
nome & grau & papel \\
\midrule
\texttt{grau1} & 1 & campo linear: todo esquema consistente é exato \\
\texttt{grau2} & 2 & o pedido original; separa malha uniforme de adaptada \\
\texttt{grau3} & 3 & separa o fluxo de segunda ordem do compacto e do MLS \\
\texttt{grau4} & 4 & separa o compacto do MLS de ordem 4 \\
\texttt{grau6} & 6 & acima de toda exatidão disponível \\
\texttt{trig}  & --- & $\sin/\cos$: não polinomial, logo há ordem a medir \\
\bottomrule
\end{tabular}
\end{center}

Os coeficientes de cada monômio são \textbf{distintos}, de propósito: assim o
valor identifica a posição, e um esquema que troque os eixos não passa por acaso.

\subsection{As normas e a ordem}

O erro é medido em duas normas, a do máximo e a $L_2$ ponderada por área:
\begin{equation}
  \|e\|_\infty = \max_c |u_c - u(\mathbf{x}_c)|,
  \qquad
  \|e\|_2 = \sqrt{\frac{\sum_c A_c\,(u_c - u(\mathbf{x}_c))^2}{\sum_c A_c}} .
  \label{eq:normas}
\end{equation}
A ponderação por área da norma $L_2$ é necessária em malha adaptada: sem ela, as
células pequenas --- que são muitas --- dominariam a média sem corresponder à
fração do domínio que ocupam.

A ordem de convergência é estimada entre níveis consecutivos de refino,
\begin{equation}
  p = \log_2 \frac{\|e\|_{h}}{\|e\|_{h/2}} .
  \label{eq:ordem}
\end{equation}

\paragraph{Onde o esquema é exato não há ordem para medir.} Isso não é uma
ressalva: é a razão pela qual as tabelas da seção \ref{sec:resultados} exibem um
travessão em vez de um número. Com erro de máquina em dois níveis consecutivos,
\eqref{eq:ordem} divide ruído por ruído e produz qualquer valor. Imprimir esse
valor sugeriria uma taxa de convergência onde existe apenas aritmética de ponto
flutuante.

\subsection{O piso do erro é o solver}

Todas as corridas usam PETSc com tolerância \num{1e-14}. O erro residual
observado onde o esquema é exato fica na faixa de \num{1e-15} a \num{1e-13}, e
ele é \textbf{resíduo do solver linear}, não do esquema. Exigir menos que isso
seria pedir ao solver algo que ele não promete com essa tolerância. Os casos da
suíte usam \num{1e-10} como limiar de ``erro de máquina'' por essa razão.

\subsection{O resíduo do operador, como oráculo independente}
\label{sec:met-residuo}

A comparação da solução discreta com a exata mede o sistema inteiro: esquema,
condições de contorno e solver juntos. Para separar o esquema do resto, há um
segundo oráculo, que não resolve nada:
\begin{equation}
  r = \max_c \frac{\bigl|(A\,\mathbf{u}^{\text{exata}})_c - b_c\bigr|}{\text{normalização}_c},
  \label{eq:residuo}
\end{equation}
isto é: monta-se o operador, aplica-se à \emph{solução exata amostrada nos centros
das células}, e compara-se com o lado direito. Se o esquema é exato para aquela
solução, $r$ é de máquina --- sem solver envolvido.

Duas propriedades fazem desse oráculo mais que uma conveniência. Primeiro, ele
distingue \textbf{consistência} de \textbf{convergência}: um operador pode ser
consistente e o sistema divergir por causa do contorno, ou vice-versa. Segundo, a
montagem é \emph{a mesma função} que o solver usa --- extraída para que os dois
caminhos compartilhem um único código. Duas montagens separadas divergiriam
exatamente onde é difícil ver: perto da parede.

O valor desse oráculo ficou patente na transição 2:1. O resíduo do operador em
malha adaptada, para grau~2, dá \num{4.5} --- contra \num{2.27e-13} na uniforme.
Esse contraste de treze ordens de grandeza é o que afirma, sem ambiguidade, que a
diferença é do \emph{esquema na transição}, não do solver nem do contorno.

\subsection{Contraste obrigatório}

Nenhuma afirmação de exatidão entra na suíte sozinha. Para cada ``o erro é de
máquina'' há um caso vizinho medindo a condição em que ele deixa de ser:

\begin{center}
\begin{tabular}{ll}
\toprule
afirmação & contraste que a acompanha \\
\midrule
grau 2 exato com fechamento de 3 pontos & grau 2 com fechamento de 2 pontos \\
grau 3 exato no MLS de ordem 3 & grau 3 no fluxo de segunda ordem \\
grau 2 exato em malha adaptada com MLS & grau 2 em malha adaptada com fluxo \\
grau 4 exato com Gauss $2\times2$ na fonte & o mesmo com a fonte no ponto médio \\
pesos próprios concordam com o herdado & concordância \emph{exatamente} nula é recusada \\
\bottomrule
\end{tabular}
\end{center}

A última linha merece explicação: um caso que afirma ``as duas implementações
concordam'' e obtém diferença exatamente zero não está comparando duas
implementações --- está comparando uma consigo mesma. Esse foi um defeito real na
comparação CPU--GPU (seção \ref{sec:est-gpu}), e desde então a concordância
perfeita é tratada como sintoma, não como sucesso.

\section{Resultados}
\label{sec:resultados}

\status{Pronta para malha uniforme e adaptada, em duas dimensões e em execução
serial. Todos os números vêm de \texttt{dados/*.dat}, gerados por
\texttt{gera-dados.sh}; nenhum foi digitado no texto. O que não foi medido ---
paralelo e três dimensões --- está na seção \ref{sec:conc-falta}, não aqui.}

Todas as corridas usam o quadrado unitário, condição de Dirichlet em todas as
paredes tomada da própria solução exata, PETSc com tolerância \num{1e-14}, e três
níveis de refino ($h = \num{1.25e-1}$, \num{6.25e-2}, \num{3.125e-2}). A malha
adaptada é uma \emph{quadtree} obtida por refino aleatório de 40\% das células com
semente fixa, balanceada em 2:1.

\subsection{Mapa dos graus de exatidão}

O resultado central, antes das tabelas: o \textbf{grau de exatidão} de cada
esquema em cada tipo de malha. É a tabela que decide qual esquema serve para quê.

\begin{center}
\begin{tabular}{lcc}
\toprule
esquema & malha uniforme & malha adaptada (2:1) \\
\midrule
fluxo, fechamento de 2 pontos & 1 & 1 \\
fluxo, fechamento de 3 pontos & \textbf{2} & 1 \\
compacto (Mehrstellen, Lele)  & \textbf{3} & \emph{recusa} \\
MLS + Gauss, ordem 3          & \textbf{3} & \textbf{3} \\
MLS + Gauss, ordem 4          & \textbf{4} & 3$^{\ast}$ \\
\bottomrule
\end{tabular}

\smallskip
\footnotesize $^{\ast}$ A ordem 4 \emph{rebaixa} em malha adaptada por falta de
pontos no suporte, e o rebaixamento é contado --- seção \ref{sec:res-rebaixo}.
\end{center}

\subsection{Malha uniforme: o esquema de fluxo}

\paragraph{Grau 2 é exato, e é o pedido original.}

\tabconv{unif-fluxo2-grau2.dat}{Grau 2, fluxo de segunda ordem, fechamento de
três pontos, malha uniforme. O erro é de máquina nas duas normas e \emph{não
melhora nem piora com o refino} --- o que é a assinatura de um esquema exato: o
que sobra é resíduo do solver, e ele não depende de $h$.}

\paragraph{O contraste do fechamento.} Com dois pontos na parede em vez de três,
o mesmo grau 2 deixa de ser exato:

\tabconv{unif-fluxo2-2pt-grau2.dat}{Grau 2, mesmo esquema, fechamento de
\textbf{dois} pontos. Erro de \num{1.16e-3} contra \num{3.91e-13} --- dez ordens
de grandeza --- e ordem $\approx 2$. Esta tabela é o que dá sentido à anterior.}

\paragraph{Grau 3 e o caso trigonométrico.}

\tabconv{unif-fluxo2-grau3.dat}{Grau 3, fluxo de segunda ordem. A ordem medida é
$\approx\num{2.95}$, acima de 2: para um campo cúbico as derivadas de quarta
ordem se anulam e o termo dominante do truncamento desaparece. Não é erro de
medida, e é por isso que o grau 3 \emph{não} serve sozinho para distinguir os
esquemas --- o caso trigonométrico serve.}

\tabconv{unif-fluxo2-trig.dat}{Solução trigonométrica, fluxo de segunda ordem.
Ordem \num{1.97} na norma do máximo. Aqui nenhuma derivada se anula, e a ordem
medida é a ordem do esquema.}

\subsection{Malha uniforme: o compacto de quarta ordem}

\tabconv{unif-lele4-grau3.dat}{Grau 3, compacto na forma de derivada segunda
implícita. Exato --- e o grau 3 é o teto, limitado pelo fantasma de quatro pontos
na parede \eqref{eq:fantasma4}, não pelo estêncil interior.}

\tabconv{unif-lele4-grau4.dat}{Grau 4, o mesmo esquema. Deixa de ser exato, e
converge em ordem \num{3.95} --- quarta ordem, como projetado. A exatidão se perde
em grau 4 e a \emph{ordem} não sofre, porque o erro do fechamento é local e também
$O(h^4)$.}

\tabconv{unif-lele4-trig.dat}{Solução trigonométrica, compacto. Ordem
\num{3.86} na norma do máximo, contra \num{1.97} do esquema de fluxo na mesma
malha e na mesma solução.}

\subsection{Malha uniforme: MLS com quadratura de Gauss}

\tabconv{unif-mls3-grau3.dat}{Grau 3, MLS de ordem 3 com Gauss. Exato, com erro
na faixa do solver.}

\tabconv{unif-mls4-grau4.dat}{Grau 4, MLS de \textbf{ordem 4} com Gauss. Exato ---
o grau de exatidão acompanha a ordem do ajuste, que é a propriedade que distingue
este esquema dos outros quatro.}

\tabconv{unif-mls4-trig.dat}{Solução trigonométrica, MLS de ordem 4. Ordem
\num{3.82} na norma do máximo; com ordem 3 o mesmo caso dá \num{2.83}. A ordem de
convergência também acompanha a ordem do ajuste.}

\subsection{O fator que destravava o esquema era a quadratura da fonte}
\label{sec:res-fonte}

Este é o achado que mais mudou o resultado, e ele é medido isoladamente trocando
\textbf{uma} coisa: a quadratura de $\int_{\Omega_c} f\,\mathrm{d}A$. A derivada
na faceta é exata nas duas corridas abaixo; a malha, a solução, o solver e a ordem
do ajuste são os mesmos.

\tabconv{unif-mls4-pm-grau4.dat}{Grau 4, MLS de ordem 4 com Gauss na faceta mas
a fonte na regra do \textbf{ponto médio}. O esquema \textbf{trava em segunda
ordem}: \num{9.58e-5} e ordem \num{1.99}, contra \num{5.91e-14} exato com Gauss
$2\times2$ --- dez ordens de grandeza de diferença.}

A regra do ponto médio é um erro $O(h^2)$ \emph{relativo} em
$\int_{\Omega_c} f\,\mathrm{d}A$. No esquema de fluxo de segunda ordem esse erro
se escondia atrás do erro do próprio fluxo, e nunca foi o termo dominante. Com o
fluxo levado à ordem alta, ele passa a ser --- e nenhum trabalho adicional sobre o
fluxo teria efeito enquanto a fonte permanecesse no ponto médio.

O sinalizador \verb|--fonte-ponto-medio| existe exatamente para produzir essa
tabela, e a cláusula D6 do contrato a exige. Sem o interruptor, ``usamos Gauss
$2\times2$ na fonte'' seria uma afirmação sem prova.

\subsection{Malha adaptada: onde os esquemas se separam}

\paragraph{O esquema de fluxo perde a exatidão do grau 2.}

\tabconv{adap-fluxo2-grau2.dat}{Grau 2, fluxo de segunda ordem, malha adaptada.
O erro é \num{1.12e-3} --- comparar com \num{3.91e-13} na malha uniforme --- e
converge em ordem \num{1.84}. A exatidão do grau 2 era propriedade da malha
uniforme.}

A causa é a quadratura: na transição, a distância entre os centros ($3h/4$) não é
a distância do centro à face ($h/2$), e o termo de curvatura não cancela. Isso
\textbf{não} é defeito, e afirmar explicitamente que a exatidão se perde é o que
impede alguém de ``consertar'' o que está correto. Na malha adaptada o esquema de
fluxo é exato apenas para campo linear.

\paragraph{O MLS com Gauss recupera a exatidão, e a estende.}

\tabconv{adap-mls3-grau2.dat}{Grau 2, MLS de ordem 3 com Gauss, \textbf{na mesma
malha adaptada} --- mesma fração, mesma semente. Exato: \num{1.52e-12} contra
\num{1.12e-3} do esquema de fluxo.}

\tabconv{adap-mls3-grau3.dat}{Grau 3, MLS de ordem 3, malha adaptada. Também
exato. O grau de exatidão em malha adaptada é o mesmo da malha uniforme, que é a
razão de existir deste esquema.}

\tabconv{adap-mls3-grau4.dat}{Grau 4, MLS de ordem 3, malha adaptada. Acima da
ordem do ajuste, logo não exato, e converge em ordem \num{3.56}.}

\tabconv{adap-mls3-trig.dat}{Solução trigonométrica, MLS de ordem 3, malha
adaptada. Ordem \num{3.08} na norma do máximo, contra \num{1.84} do esquema de
fluxo na mesma malha.}

\subsection{A ordem 4 não sobrevive à malha adaptada}
\label{sec:res-rebaixo}

Uma medida negativa, e ela é um resultado. Em malha adaptada com ordem 4,
\textbf{8 dos \num{5172}} ajustes MLS não sustentam a ordem pedida. A causa é o
\emph{suporte}, não o condicionamento: o crescimento por anéis entrega de 12 a 16
pontos, contra os 15 monômios que a ordem 4 exige. Em malha uniforme o mesmo
crescimento entrega suporte regular e não há rebaixamento algum.

A consequência é que a exatidão de grau 4 se perde na malha adaptada --- e deve
se perder: oito pontos de quadratura usam ordem menor que a declarada. A ordem
utilizável em malha adaptada é~3, e levá-la a~4 é um trabalho de \emph{seleção de
suporte} --- mais pontos, escolhidos por qualidade geométrica em vez de por anéis
--- não de quadratura nem de fechamento.

O que torna essa medida possível é o contador de rebaixamento descrito na seção
\ref{sec:aprox-mls}. Antes dele, o mesmo caso produzia o defeito relatado na
seção \ref{sec:verif-defeitos}.

\subsection{Conservação na faceta de transição}

A cláusula D7 compara, nas sub-facetas de transição 2:1, os pontos de quadratura
e os estênceis de derivada vistos dos \emph{dois} lados --- com o oráculo
recalculando a geometria por conta própria. Em oito sub-facetas de transição, a
diferença máxima é \textbf{\num{0.0e+00}} no ponto e \textbf{\num{0.0e+00}} no
peso: identidade exata, não concordância numérica.

Isso é consequência do desenho, não de um acerto: o suporte depende do ponto, não
de quem pergunta. E merece oráculo próprio porque um esquema não conservativo
\emph{ainda converge} --- a propriedade é invisível para toda a bateria de
tabelas acima.

\section{A suíte, e o que ela já pegou}
\label{sec:verificacao}

\status{Pronta. O estado relatado é o da execução de 12 de outubro de 2026: 141
casos, nenhum falhando, contrato 69/69. Os defeitos da seção
\ref{sec:verif-defeitos} são todos reais, e cada um está nomeado com o oráculo que
o encontrou.}

\subsection{Contrato lido pelo arreio}

A suíte é organizada como um \textbf{contrato}: um cabeçalho com cláusulas
numeradas, cada uma declarando uma propriedade que o sistema afirma, explicando
por que ela merece ser afirmada, e nomeando o caso de teste que a verifica. O
arreio de execução \emph{lê as cláusulas do próprio cabeçalho} e exige um caso
passando para cada uma --- cláusula sem teste sai como \textbf{falha}, não como
omissão.

Isso não é formalidade. Em um dos passos deste trabalho, um refator removeu o caso
que cobria uma cláusula, e o arreio reportou ``caso ausente'' nomeando a cláusula
--- o que uma suíte que apenas conta testes passando não faria.

O estado atual, em duas dimensões:

\begin{center}
\begin{tabular}{llr}
\toprule
grupo & assunto & cláusulas \\
\midrule
N & mínimos quadrados móveis: valor, derivada, recusas & 7 \\
B, F & floresta, blocos, suporte & 4 \\
M, R & malha, numeração, refino & 13 \\
G, A, L & GPU, arranjo deslocado, lote de estênceis & 15 \\
E, C & estêncil, condições de contorno por região & 12 \\
P, Q, T & Poisson, compacto, malha adaptada & 11 \\
D & derivada de ordem alta e quadratura de Gauss & 7 \\
\midrule
\multicolumn{2}{l}{\textbf{total}} & \textbf{69} \\
\bottomrule
\end{tabular}
\end{center}

Quatro guardas no arreio merecem menção, porque cada uma existe por um modo de
falha observado: ela nomeia o prefixo de cláusula sem cobertura em vez de apenas
somar; faz \verb|make clean| por configuração, para que um binário compilado com
outro sinalizador não seja medido como se fosse deste; exige a linha de resumo do
executável, para que um processo que morreu antes de terminar não conte como
aprovado; e executa tudo por \verb|mpirun| mesmo com um processo, para que o
caminho de execução seja o mesmo do paralelo.

\subsection{Quatro defeitos que só o contraste encontrou}
\label{sec:verif-defeitos}

Esta subseção existe porque a disciplina de oráculo é barata de descrever e difícil
de levar a sério, e o argumento mais convincente a seu favor é a lista do que ela
pegou neste trabalho.

\paragraph{1. O estêncil medido onde o ajuste é degenerado.} Os primeiros testes
de estêncil e de contorno mediam o ajuste \emph{no centro da célula}. Ali o MLS é
degenerado por construção --- peso $1$ na própria célula, lado direito
$\approx 0$ --- de modo que o fechamento de contorno \textbf{nunca era
exercitado}. Os casos passavam. O que denunciou foi o contraste: ele dava apenas
\num{1e-10} quando devia dar ordens de grandeza. Com alvos deslocados do centro, o
contraste passou a \num{1.9e-1}, e uma cláusula nova agora \emph{afirma} a
degeneração no centro, para que ninguém a redescubra como surpresa.

\paragraph{2. A correção transversal com sinal dependente de quem é grosso.} Na
transição 2:1, fazer o sinal da correção transversal depender de qual lado é o
grosso \textbf{dobrou} o erro: \num{7.8e-2} contra \num{3.9e-2}. A fórmula correta
é independente disso. O sintoma --- erro exatamente duas vezes maior --- foi o que
identificou a causa.

\paragraph{3. O \emph{nan} silencioso do ajuste de ordem 4.} Este é o mais grave,
e é o que motivou o contador de rebaixamento. A rotina de eliminação recusava
pivô somente abaixo de \num{1e-300} --- isto é, apenas a singularidade exata --- e,
ao recusar, devolvia um \textbf{vetor de zeros}. Em malha uniforme isso nunca
ocorria. Em malha adaptada com ordem 4, ocorria em 8 de \num{5172} ajustes, e o
efeito era uma linha da matriz silenciosamente errada:

\begin{center}
\begin{tabular}{lr}
\toprule
nível & $\|e\|_\infty$ (antes do conserto) \\
\midrule
3 & \num{2.36e-10} \\
4 & \textbf{\num{3.32e-1}} \\
5 & \num{2.35e-10} \\
\bottomrule
\end{tabular}
\end{center}

Um nível com erro de \num{3.3e-1} cercado por \num{1e-10} nos vizinhos. O conserto
tem duas partes --- a base adimensionalizada e o rebaixamento contado --- e ambas
estão descritas na seção \ref{sec:aprox-mls}.

\paragraph{4. O ponteiro nulo no ramo que não o usa.} Ao estender a montagem para
o esquema novo, uma edição por substituição de texto inseriu a leitura dos
contadores de MLS \emph{também} no ramo do esquema compacto, onde o objeto
correspondente é um ponteiro nulo. O compilador não acusou, o esquema novo
funcionava, e o compacto --- que passava antes --- quebrava com violação de
segmentação. O experimento que localizou isso em um passo foi trocar a
\emph{ordem} dos casos no arreio: a falha permaneceu no mesmo ponto lógico, o que
descartou de uma vez a hipótese de acumulação de estado e apontou para o caso.

\subsection{O padrão, destilado}

Os quatro têm a mesma forma: \textbf{a afirmação estava certa e o oráculo não
discriminava}. Em três dos quatro, o que denunciou não foi um caso falhando --- foi
um \emph{contraste insuficiente}, isto é, um número que devia ser grande e era
pequeno. Daí a regra de que nenhuma afirmação de exatidão entra na suíte sem o par
que mede onde ela falha.

\section{A execução paralela}
\label{sec:paralelo}

\status{Pronta. As cláusulas X1--X9 e H18--H19 cobrem a correção, e o arreio as
executa em np = 1, 2, 3 e 4; o custo da franja está na seção
\ref{sec:par-custo} e a escalabilidade do sistema inteiro na seção
\ref{sec:par-escala}, esta última limitada a oito núcleos físicos.}

Esta seção é separada das anteriores por um motivo metodológico: tudo o que ela
descreve é invisível em execução serial, e \textbf{cinco} defeitos reais do
sistema viviam exatamente nessa invisibilidade. Relatá-los é o conteúdo mais
transferível do trabalho.

Os três primeiros apareceram em duas dimensões; os dois últimos, em três --- e o
quarto é o segundo voltando por baixo do seu próprio conserto, o que é a razão
de valer a pena relatá-lo.

\subsection{A franja de $K$ camadas}

A camada de fantasmas nativa da biblioteca tem \textbf{uma} camada, e uma camada
não basta. O ponto de quadratura de Gauss está \emph{sobre} a faceta, logo a
célula que o contém pode ser de franja --- e a partir dela o suporte do ajuste MLS
ainda precisa abrir os anéis inteiros. Medido com uma camada:

\begin{center}
\texttt{suporte de 1 pontos em 6 aneis, minimo 10 para ordem 3}
\end{center}

As camadas adicionais vêm de \textbf{rodadas de troca}: cada processo publica, por
folha sua, a vizinhança que já conhece, e cada rodada estende o alcance em um
anel. A regra é $K = \text{anéis} + 1$, e a diferença de um foi \emph{medida}, não
deduzida --- com $K$ igual ao número de anéis, np=2 truncou 32 suportes e np=3
falhou.

\paragraph{A alternativa rejeitada.} Crescer o suporte a partir da célula local
adjacente ao ponto resolveria o problema e destruiria a conservação: o suporte
passaria a depender de \emph{quem pergunta}, e os dois fluxos pela mesma faceta
deixariam de se cancelar --- perdendo exatamente a propriedade que a seção
\ref{sec:aprox-mlsgauss} obtém por construção.

\subsection{Defeito 1: a franja ordenada por id}

O mais caro, e o mais instrutivo porque \textbf{o comentário advertindo contra ele
estava escrito na mesma função}.

A documentação da biblioteca é explícita: a consulta de vizinhos numera
$0 \ldots n_{\text{loc}}-1$ para as folhas locais e
$n_{\text{loc}} \ldots n_{\text{loc}}+n_{\text{fan}}-1$ para as da primeira
camada de franja. O achatamento da tabela interna ordenou \emph{toda} a franja por
id --- porque id é determinístico e parecia a escolha limpa. O efeito: o índice
devolvido para o terceiro fantasma passou a designar outra folha, e a consulta
entregava um id que \textbf{existe}, só que de outra célula.

\begin{center}
\begin{tabular}{lrr}
\toprule
caso & np=1 & np=2 \\
\midrule
\texttt{grau1} adaptado, fluxo de 2ª ordem & \num{3.97e-13} & \textbf{\num{1.70e-1}} \\
\bottomrule
\end{tabular}
\end{center}

O esquema é \textbf{exato} para campo linear, logo o número de processos não podia
mudar nada --- e é essa impossibilidade que torna o caso um oráculo. Somente as
camadas 2 e acima podem ser ordenadas por id: elas não têm índice da biblioteca.

\subsection{Defeito 2: face de franja indistinguível de contorno}

Uma face sem vizinhos na lista significa \textbf{duas} coisas: contorno do
domínio, ou ``o dono desta folha me contou sobre as outras faces e não sobre
esta''. Com um indicador booleano por folha as duas se confundiam, e a segunda era
tratada como a primeira: o esquema aplicava condição de contorno \textbf{no meio
do domínio}, ao longo de toda a fronteira de partição, em silêncio.

O conserto é uma \textbf{máscara de faces conhecidas}, um bit por face, publicada
junto com os vizinhos na troca.

\subsection{Defeito 3: o retrocesso silencioso}

O gradiente transversal da transição 2:1 precisa do nível do vizinho transversal
da célula grossa. A primeira versão o buscava na tabela e, quando faltava, caía
num valor de reserva --- o nível da própria célula grossa. Isso dava espaçamento
errado, distância errada, gradiente errado.

Duas lições, e a segunda custou mais que a primeira. A primeira: \textbf{o valor
de reserva plausível foi pior que uma recusa}. A segunda: a busca pela causa foi
para o lugar errado. Tentei mais camadas de franja, até quatro, e as ausências não
iam a zero --- porque não eram camadas: era a \textbf{carga da troca saturando},
já que em malha adaptada a vizinhança de uma célula grossa tem muitos vizinhos
finos. O conserto foi fazer o nível \emph{viajar junto} com o id do vizinho, o que
economiza uma camada inteira, e \textbf{contar} a saturação em vez de cortá-la
calada.

\subsection{Defeito 4: a máscara que mentia}

O defeito 2 foi consertado com uma máscara de faces conhecidas, um bit por face,
publicada junto com os vizinhos. Em três dimensões a máscara \textbf{mentia}, e
o defeito 2 voltou por baixo do seu próprio conserto.

A máscara era escrita \emph{antes} da carga útil, e a carga era truncada por
capacidade. O campo da máscara entrou no bloco publicado sem que a conta de
tamanho crescesse: cabiam $\mathtt{max\_viz}-1$ vizinhos onde a folha tem
$\mathtt{max\_viz}$. Numa malha uniforme, onde cada folha tem seis vizinhos, o
\textbf{sexto} ficava para trás --- e o sexto é a face 5, a última do laço. A
folha chegava afirmando conhecer as seis faces com a lista da sexta ausente, e
\emph{lista vazia com bit ligado é exatamente a assinatura de uma parede do
domínio}. O fechamento aplicava Dirichlet no meio do domínio, de novo.

O conserto tem duas partes, e a segunda importa mais que a primeira. A capacidade
foi corrigida; e a máscara passou a ser escrita \textbf{depois} da carga,
afirmando apenas as faces que couberam inteiras. \textbf{Uma face parcial é pior
que uma face desconhecida}: desconhecida, o fechamento se recusa a aplicar
condição e o esquema para; parcial, ele aplica a condição errada e o resultado
sai plausível. O mesmo endurecimento foi feito em duas dimensões, onde a carga
cabe e portanto o defeito nunca disparou --- a diferença entre ``não dispara'' e
``não pode disparar'' é uma conta que alguém tem de refazer a cada mudança de
tamanho, e é justamente o tipo de conta que não se refaz.

\subsection{Defeito 5: uma rodada de troca só alcança vizinho de face}

A franja nativa do t8code é \texttt{T8\_GHOST\_FACES}, e a documentação de
\texttt{t8\_forest\_set\_ghost} registra que \emph{só} esse tipo é suportado.
Logo o processo diagonalmente oposto numa quina não é vizinho de face de ninguém
que fale comigo, e uma rodada de publicação nunca traz a célula dele. Em oito
processos isso aparecia como 192 suportes truncados, todos nas células que tocam
o centro do domínio --- onde as oito partições se encontram.

\textbf{Aumentar as camadas não resolvia}, e foi isso que confundiu o
diagnóstico: com 2, 3 e 4 camadas o truncamento ficava preso em 192. A informação
não precisava ir mais \emph{longe}; precisava ir por outro \emph{caminho}.

A fase de publicação passa a \textbf{iterar até cobrir}, com parada medida e
coletiva: quantos pontos da caixa ainda faltam, reduzido entre todos os
processos. Cada rodada avança um salto de posse, e o número de rodadas sai da
partição, não de uma constante --- medido no nível 3: np=2 pede 1 rodada, np=4
pede 2, np=8 pede 3. Um número fixo estaria certo para esta partição e errado
para a próxima.

O mesmo trabalho expôs uma terceira peça. O esquema pedia $r+1$ camadas e precisa
de $r+2$: a caixa do suporte não é centrada na folha local, e sim na célula que
contém o \textbf{ponto de quadratura}, que fica sobre uma faceta e pode cair na
célula vizinha. Com $r+1$ o esquema sobrevivia por acidente --- a publicação em
torno dos fantasmas cobria o excedente sem que nada afirmasse isso --- e o
acidente sumiu quando a publicação passou a parar no critério declarado.

\subsection{O oráculo que faltava, e o que ele exigiu}
\label{sec:par-oraculo}

Os defeitos 4 e 5 ficaram invisíveis por um motivo que não é sobre eles: o
arquivo de teste do Poisson tridimensional \textbf{só rodava em um processo}, e
em um processo não há franja. A suíte ficou verde com o paralelo errado.

O primeiro oráculo que tentamos não discriminava. O resíduo direto do operador,
$\max|A\,u_{\text{exata}} - b|$, é um \emph{máximo}, dominado pelas linhas de
parede: ficou idêntico em \num{1.1591} nos quatro números de processos enquanto
o interior divergia.

O que achou foi a \textbf{soma da linha}. Um laplaciano aniquila constantes,
logo a linha de uma célula interior soma zero --- e isso é propriedade do
operador, verificável em qualquer número de processos, sem referência serial. Em
série somava \num{1e-17}; em paralelo, \num{-1.907e-2} em quatro linhas. Soma
diferente de zero significa peso \emph{saindo} da linha, indo para o lado direito
como condição de contorno, que é exatamente o sintoma do defeito 4.

O caso novo afirma as duas coisas que não precisam de serial --- soma de linha
interior nula e zero suportes truncados --- e roda em np=1, 2 e 4. Verificamos
que ele \textbf{discrimina}: com os defeitos reintroduzidos, passa em np=1 e
reprova em np=2.

\reserva{o defeito ficou mais um dia escondido por uma causa que não é de
numérica: o \texttt{Makefile} não rastreava dependência de cabeçalho. A função
que calcula as camadas é \emph{inline} no cabeçalho; mudou de $r+1$ para $r+2$ e
o objeto do Poisson continuou pedindo $r+1$. A medida dizia que o conserto não
tinha efeito quando o conserto não estava no binário --- não há erro, há uma
resposta errada com cara de medida.}

\subsection{O que está verificado}

Mesmo problema, mesmo resultado, variando apenas o número de processos:

\begin{center}
\begin{tabular}{lrrr}
\toprule
caso & np=1 & np=2 & np=3 \\
\midrule
\texttt{grau2} uniforme, fluxo 2ª ordem        & \num{2.99e-13} & \num{2.75e-13} & \num{1.15e-13} \\
\texttt{grau2} \textbf{adaptado}, fluxo 2ª ordem & \num{3.9960e-3} & \num{3.9960e-3} & \num{3.9960e-3} \\
\texttt{grau3} uniforme, MLS+Gauss             & \num{7.99e-15} & \num{2.18e-14} & \num{3.81e-14} \\
\texttt{grau2} \textbf{adaptado}, MLS+Gauss      & \num{7.02e-13} & \num{3.54e-13} & \num{1.55e-13} \\
\texttt{grau4} uniforme, MLS ordem 4           & \num{2.52e-14} & \num{2.11e-13} & \num{2.51e-14} \\
\bottomrule
\end{tabular}
\end{center}

\textbf{A segunda linha é o oráculo mais forte da tabela.} Um esquema que
\emph{não} é exato dando erro \emph{idêntico} nos três casos significa matriz
bit-a-bit igual --- nenhuma das outras linhas afirma tanto, porque onde o esquema é
exato o que varia é resíduo do solver linear, que legitimamente depende da
partição (o precondicionamento depende dela).

Também verificado: as normas e o resíduo reduzem entre processos; o resíduo
percorre a linha inteira --- antes descartava as colunas fora da faixa local, e
saía \emph{menor} que o verdadeiro, isto é, o oráculo mentindo a favor do código;
a repartição preserva a contagem global, a área, os níveis e o número de células
em transição, e refaz a franja; e a correspondência id\,$\to$\,geometria não
depende da partição.

\subsection{O custo da franja}
\label{sec:par-custo}

Tempo absoluto de franja não é resultado: vinte milissegundos podem ser
desprezíveis ou dominantes. O que interpreta é a fração do tempo de resolver o
sistema.

\begin{center}
\begin{tabular}{rrrrrr}
\toprule
nível & células & np & franja (3 camadas) & \% do fluxo & \% do MLS \\
\midrule
6 & \num{4096}  & 2 & \SI{19.0}{\milli\second} & 33\,\% & 7,2\,\% \\
6 & \num{4096}  & 8 & \SI{10.9}{\milli\second} & 30\,\% & 12\,\% \\
7 & \num{16384} & 2 & \SI{93.1}{\milli\second} & 33\,\% & \textbf{3,5\,\%} \\
7 & \num{16384} & 8 & \SI{32.3}{\milli\second} & 35\,\% & \textbf{9,6\,\%} \\
\bottomrule
\end{tabular}
\end{center}

\textbf{O resultado que decide a questão está na coluna do MLS, na passagem do
nível 6 para o 7: a fração CAI com o tamanho do problema} --- de 7,2\,\% para
3,5\,\% em np=2. A franja cresce com o número de folhas; o solver do MLS cresce
mais rápido, porque sua matriz tem da ordem de 300 não-zeros por linha. No regime
que importa --- muitas células por processo --- a franja é ruído.

E a comparação é mais favorável do que a tabela mostra, por algo que ela não
captura: a franja é remontada quando a \emph{malha muda}, e o sistema é resolvido
a \emph{cada passo}.

A medida produziu duas otimizações que não estavam previstas. A primeira:
\texttt{exige\_camadas} em um processo só refazia todo o trabalho local --- 
\SI{22}{\milli\second} numa malha de \num{4096} células, 26\,\% do tempo de
resolver --- para não trocar dado nenhum; sem fantasmas não há o que refazer, e
hoje a chamada apenas registra o número de camadas. A segunda: o buffer da troca
era dimensionado por $(2K+1)^2$ em \emph{todas} as rodadas, embora a rodada $r$
publique apenas a vizinhança de raio $r$; dimensioná-lo pela rodada reduziu o
volume publicado em \num{2.1} vezes.

\subsection{A escalabilidade, e o superlinear que não era}
\label{sec:par-escala}

Uma versão anterior deste texto registrava, como ressalva, um ganho aparente de
\num{29.8}$\times$ de np=1 para np=8 no esquema MLS --- \emph{superlinear} em
oito processos, portanto impossível. A hipótese declarada na época era cache ou
mudança no número de iterações. \textbf{Era nenhuma das duas}: a montagem do
operador era \textbf{quadrática}. A busca geométrica que converte ponto em
índice era uma varredura linear, chamada uma vez por linha; com índice espacial,
o mesmo problema caiu de \SI{8.74}{\second} para \SI{0.679}{\second} e o ganho
aparente virou \num{5.2}$\times$ --- sublinear, como tem de ser.

A lição não é sobre o algoritmo: é que \textbf{um ganho impossível não é um
bônus, é um sintoma}, e o que ele indicava era um defeito de complexidade no
caminho mais quente do código, disfarçado de boa notícia.

Com isso resolvido, a medida tridimensional. Oito núcleos \textbf{físicos} nesta
máquina --- os dezesseis que o sistema relata são fios, e medir acima de oito
mede contenção de fio, não escalabilidade:

\begin{center}
\begin{tabular}{lrrrrrrr}
\toprule
esquema & células & np & franja & montagem & resolução & total & ganho \\
\midrule
fluxo 2ª ordem & \num{262144} & 1 & --- & \SI{5.40}{\second} & \SI{2.60}{\second} & \SI{8.00}{\second} & --- \\
               &              & 2 & --- & \SI{3.54}{\second} & \SI{1.83}{\second} & \SI{5.37}{\second} & \num{1.49}$\times$ \\
               &              & 4 & --- & \SI{1.60}{\second} & \SI{0.96}{\second} & \SI{2.57}{\second} & \num{3.12}$\times$ \\
               &              & 8 & --- & \SI{0.93}{\second} & \SI{0.69}{\second} & \SI{1.62}{\second} & \num{4.93}$\times$ \\
\midrule
MLS+Gauss      & \num{32768}  & 1 & --- & \SI{6.30}{\second} & \SI{0.79}{\second} & \SI{7.09}{\second} & --- \\
               &              & 2 & \SI{0.12}{\second} & \SI{3.27}{\second} & \SI{0.48}{\second} & \SI{3.87}{\second} & \num{1.83}$\times$ \\
               &              & 4 & \SI{0.09}{\second} & \SI{1.73}{\second} & \SI{0.20}{\second} & \SI{2.02}{\second} & \num{3.52}$\times$ \\
               &              & 8 & \SI{0.07}{\second} & \SI{0.90}{\second} & \SI{0.12}{\second} & \SI{1.09}{\second} & \num{6.52}$\times$ \\
\bottomrule
\end{tabular}
\end{center}

O esquema de fluxo não tem coluna de franja porque pede uma camada só e não
entra na fase iterativa --- a célula vazia é informação, não omissão.

\textbf{Nos dois esquemas o erro em np=8 é igual ao de np=1 com desvio relativo
\num{0.00e0}}, e é essa linha que autoriza a leitura das outras: ganho sobre
resultado diferente não é ganho. As três fases são cronometradas separadas de
propósito --- somadas, uma fase que não comunica dominando em np=1 faz o total
parecer escalar bem sem que a parte difícil escale.

Ressalva declarada: o número de iterações do solver muda com np (60 para 91 no
fluxo, 25 para 35 no MLS), logo parte do que se vê é o precondicionador por
blocos, não paralelismo.

\subsection{A franja iterativa, que era o limite}
\label{sec:par-otimiza}

Corrigir a resposta custou escala. A fase iterativa que consertou a quina de
oito processos custava \SI{0.74}{\second} em np=8 --- \SI{38}{\percent} do
tempo da corrida --- e não diminuía com o número de processos. Era o teto do
esquema MLS, e os números da tabela acima já são os de depois de removê-lo.

A medida por fase, antes de otimizar, desloca a atenção inteira:

\begin{center}
\begin{tabular}{lrrlrr}
\toprule
fase & antes & & fase & antes & \\
\midrule
publica & \SI{0.296}{\second} & 44\,\% & troca & \SI{0.062}{\second} & 9\,\% \\
zera    & \SI{0.138}{\second} & 21\,\% & funde & \SI{0.015}{\second} & 2\,\% \\
conta   & \SI{0.080}{\second} & 12\,\% & resto & \SI{0.075}{\second} & 11\,\% \\
\bottomrule
\end{tabular}
\end{center}

\textbf{A troca era \SI{9}{\percent}.} O custo não era comunicar, era
\emph{preparar} o que se comunica --- e três dos quatro maiores itens eram
trabalho jogado fora:

\begin{itemize}
\item publicava-se a caixa de \emph{toda} folha local, e a troca só leva o
      bloco das que são fantasma de alguém: dois terços das folhas eram
      preenchidas, apagadas e ignoradas. O critério de quem publica é
      \textbf{exato}, não heurística --- com \texttt{T8\_GHOST\_FACES}, sou
      fantasma de $P$ exatamente quando algum vizinho meu de face é fantasma
      para mim;
\item eram \emph{dois} percursos idênticos sobre a mesma caixa, um para o
      critério de parada e outro para a carga, e o percurso é o que custa:
      $(2R+1)^3$ buscas geométricas por folha. Contar durante o preenchimento
      não muda o que se mede --- os dois liam o mesmo estado;
\item \textbf{folha coberta continua coberta}: a tabela só cresce ao longo das
      rodadas, logo percorrer de novo uma folha já inteira confirma o que já se
      sabia. Quem publica percorre sempre, porque o que ela \emph{sabe} cresce;
\item a desduplicação alocava --- um \texttt{std::set} construído e destruído
      por folha e por rodada;
\item e a fase inicial era refeita a cada chamada: o construtor monta a franja
      de uma camada e \texttt{exige\_camadas} refazia tudo, inclusive as seis
      consultas de vizinhança ao t8code por folha local (98 mil chamadas em
      np=2, \SI{55}{\percent} do tempo ali). O resultado delas já está nos
      arranjos achatados, e eles são até mais ricos.
\end{itemize}

\reserva{uma otimização foi tentada e \textbf{revertida}, e fica registrada
porque o raciocínio que a motivou é convincente e está errado. A fusão é a única
leitora das vagas não usadas e só lê o primeiro campo de cada uma, logo marcar
só esse campo deveria substituir o apagamento em bloco. Ficou \emph{mais lento}
--- \SI{0.062}{\second} para \SI{0.088}{\second} ---, porque a escrita espaçada
de um campo a cada \texttt{CAMPOS} doubles toca exatamente as mesmas linhas de
cache que o apagamento sequencial, sem economizar tráfego nenhum, e ainda perde
a vetorização. \textbf{Economizar escritas não é o mesmo que economizar memória
tocada.}}

O resultado, no mesmo caso:

\begin{center}
\begin{tabular}{lrrr}
\toprule
np & antes & depois & \\
\midrule
2 & \SI{0.805}{\second} & \SI{0.148}{\second} & \num{5.4}$\times$ \\
4 & --- & \SI{0.154}{\second} & \\
8 & \SI{0.666}{\second} & \SI{0.186}{\second} & \num{3.6}$\times$ \\
\bottomrule
\end{tabular}
\end{center}

Em np=8 a troca passou de \SI{9}{\percent} a \SI{33}{\percent} \emph{sem mudar
de valor absoluto}: o que sobrou da franja era comunicação --- e a comunicação
virou, então, o alvo seguinte.

\subsection{Uma mensagem por destinatário}
\label{sec:par-mensagem}

A troca era \texttt{t8\_forest\_ghost\_exchange\_data}: um bloco de tamanho
fixo por folha, com a caixa de raio $R$ em torno dela, o mesmo conteúdo indo
para todos. Correta, e cara por duas razões independentes.

\textbf{A célula ia repetida.} As caixas de folhas vizinhas se sobrepõem quase
inteiramente, e uma célula aparecia em até $(2R+1)^3$ blocos --- 125 cópias da
mesma geometria na mesma mensagem. \textbf{E ia o que o destinatário já tinha}:
metade da caixa de uma folha de interface é formada por células \emph{locais}
de quem recebe, e a versão local é a autoritativa, de modo que chegavam para ser
descartadas.

Agora cada célula vai uma vez por destinatário, e não vai para quem já a tem
como local --- o dono sai do id, porque a partição é contígua em id e os inícios
são trocados uma vez. O que torna a substituição possível sem mexer no resto é
que \textbf{quem recebe não perde nada}: a fusão já percorria os blocos célula a
célula e ignorava de qual folha tinham vindo. A estrutura de caixa existia só no
transporte.

A troca é ponto a ponto e só com os vizinhos --- um \texttt{Alltoall} custaria
$O(P)$ por rodada em cada processo mesmo quando cada um fala com três. A
vizinhança é \textbf{simétrica}, já que adjacência de face é simétrica, logo
quem me envia é quem eu envio e não há o que descobrir.

\begin{center}
\begin{tabular}{lrrr}
\toprule
fase & antes & depois & \\
\midrule
publica & \SI{0.062}{\second} & \SI{0.049}{\second} & \\
troca   & \SI{0.062}{\second} & \SI{0.004}{\second} & \num{14}$\times$ \\
apaga   & \SI{0.027}{\second} & ---                  & não há buffer fixo \\
funde   & \SI{0.015}{\second} & \SI{0.003}{\second} & \\
\midrule
total   & \SI{0.186}{\second} & \SI{0.074}{\second} & \\
\bottomrule
\end{tabular}
\end{center}

E a franja passou a \textbf{diminuir} com o número de processos, que é o que
não acontecia: \SI{0.121}{\second} em np=2, \SI{0.090}{\second} em np=4,
\SI{0.074}{\second} em np=8. O volume publicado em np=8 é \SI{11.67}{\mega\byte}
somados sobre os processos; a forma anterior publicava, pela sua própria
contabilidade, \SI{682}{\mega\byte} no mesmo caso --- cinquenta e oito vezes
mais.

\begin{center}
\begin{tabular}{lrrr}
\toprule
esquema MLS, np=8 & original & 1ª rodada & 2ª rodada \\
\midrule
ganho            & \num{3.70}$\times$ & \num{5.87}$\times$ & \num{6.52}$\times$ \\
eficiência       & \SI{46}{\percent}  & \SI{73}{\percent}  & \SI{82}{\percent} \\
franja no total  & \SI{38}{\percent}  & \SI{15}{\percent}  & \SI{6.6}{\percent} \\
\bottomrule
\end{tabular}
\end{center}

E a verificação que acompanha cada passo: o erro em np=8 continua idêntico ao de
np=1, desvio relativo \num{0.00e0}. \textbf{Otimização que muda o resultado não
é otimização.}

\paragraph{np=3 está na suíte por motivo medido.} Foi ele, e não np=2, que expôs a
célula grossa da transição 2:1 na borda da franja. Número par de processos parte o
quadrado ao meio e produz fronteiras alinhadas com a malha; ímpar não.

\section{Conclusões}
\label{sec:conclusoes}

\status{Pronta. Cada conclusão abaixo remete à seção que a mede; as pendências
estão separadas em \ref{sec:conc-falta} e não misturadas com os resultados.}

\subsection{O que foi estabelecido}

\paragraph{A floresta como estrutura é viável, e a herança é real.} O invólucro
de solvers lineares e o aproximador por mínimos quadrados do sistema antecessor
do sistema antecessor~\cite{sousa2019} são genuinamente independentes de malha, compilam sem modificação sob o padrão do
t8code, e produzem os mesmos pesos da implementação própria em \num{1.249e-16}
(seção \ref{sec:aprox-mls}). O que não é herdável é a propriedade da estrutura ---
e é nisso que o sistema novo difere.

\paragraph{O grau de exatidão é a propriedade que separa os esquemas.} Não a
ordem de convergência. Os cinco esquemas medidos têm ordens que se parecem em
alguns casos e graus de exatidão que não se parecem em nenhum: o esquema de fluxo
é exato em grau 2 na malha uniforme e em grau 1 na adaptada; o compacto é exato em
grau 3 e recusa a malha adaptada; o MLS com Gauss é exato até o grau da ordem do
ajuste em ambas (seção \ref{sec:resultados}).

\paragraph{A quadratura do termo de fonte era o limite, e não o fluxo.} Com a
derivada na faceta já exata, a regra do ponto médio no volume mantinha o esquema
inteiro em segunda ordem --- \num{9.58e-5} e ordem \num{1.99}, contra exato com
Gauss $2\times2$. Nenhum trabalho adicional sobre o fluxo teria efeito enquanto a
fonte permanecesse ali (seção \ref{sec:res-fonte}). Esse é o resultado deste
trabalho que menos se antecipa da análise e mais se impõe da medida.

\paragraph{Conservação na transição 2:1 sai do desenho.} Fazer o suporte do
estêncil depender apenas do ponto de quadratura, e nunca da célula que pergunta,
dá fluxo idêntico visto dos dois lados --- diferença \num{0.0e+00}, identidade
exata --- inclusive onde a célula grossa parte a faceta e a fina não parte a sua.
A propriedade precisa de oráculo próprio porque um esquema não conservativo ainda
converge.

\paragraph{Medida negativa também é resultado.} A ordem 4 do MLS não sobrevive à
malha adaptada: 8 de \num{5172} ajustes rebaixam, por falta de pontos no suporte
--- de 12 a 16 contra os 15 monômios exigidos. A ordem utilizável em malha
adaptada é 3, e a via para 4 é seleção de suporte, não quadratura (seção
\ref{sec:res-rebaixo}).

\subsection{O que este trabalho não traz}
\label{sec:conc-falta}

Declarado aqui, e não diluído entre os resultados:

\begin{itemize}
\item \textbf{a medida de custo do paralelo.} A execução paralela em si está
      feita e verificada (seção \ref{sec:paralelo}); o que falta é o
      \emph{custo}: a escalabilidade de \num{1.6}$\times$ em np=4 citada na
      seção \ref{sec:estrategias} foi medida \emph{antes} da franja de $K$
      camadas, e portanto não vale mais;
\item \textbf{três dimensões.} A base de código tridimensional não existe ainda,
      por decisão de projeto (seção \ref{sec:estrategias}). A medida de
      escalabilidade que importa --- muito trabalho por processo --- está nesse
      regime;
\item \textbf{Robin dentro do operador de Poisson.} A relação afim
      \eqref{eq:bc-afim} trata os três tipos, e o estêncil a usa; o operador de
      Poisson monta Dirichlet e Neumann. Montar Robin como se fosse Neumann daria
      resultado plausível e errado, e por isso ele recusa em vez de aproximar;
\item \textbf{fechamento de parede de ordem 5}, que elevaria a exatidão do
      compacto de grau 3 para grau 5. O limite medido é o fechamento, não o
      interior (seção \ref{sec:resultados}), de modo que esse é o passo natural
      dessa família.
\end{itemize}

\subsection{O que transferimos de método}

Fechamos com o que consideramos o resultado mais transferível, e ele não é
numérico. Em todos os defeitos deste trabalho (seção \ref{sec:verif-defeitos}), o
que separou causa de sintoma foi experimento controlado com \textbf{uma} variável
--- mesma malha, mesmo binário, mesmo solver, mudando só o que se quer isolar ---
e oráculo com \textbf{contraste}. Afirmações de exatidão sem o par que mede onde
elas falham não distinguem um esquema correto de um oráculo cego: em três dos
quatro defeitos, o que denunciou não foi um caso falhando, e sim um contraste que
devia ser grande e era pequeno.


\bibliographystyle{plain}
\bibliography{referencias}

\end{document}
